\documentclass[11pt,a4paper]{article}

\usepackage[a4paper,margin=1in]{geometry}
\usepackage{amsmath,amssymb,mathtools}
\usepackage{array,booktabs,longtable}
\usepackage{enumitem}
\usepackage{hyperref}
\usepackage{xcolor}
\usepackage[T1]{fontenc}
\usepackage[utf8]{inputenc}
\usepackage{lmodern}

\hypersetup{colorlinks=true,linkcolor=blue!50!black,urlcolor=blue!50!black}
\setlength{\parindent}{0pt}
\setlength{\parskip}{0.65em}
\renewcommand{\arraystretch}{1.25}

\newcommand{\VFS}{\mathrm{VFS}}
\newcommand{\DA}{\mathcal D_A}
\newcommand{\LV}{\mathcal L_{\rm VFS}}
\newcommand{\Rcal}{\mathcal R}
\newcommand{\Grec}{\mathcal G_{\rm recepta}}
\newcommand{\Igate}{I_{\rm gate}}
\newcommand{\Sq}{S_q}
\newcommand{\Om}{\Omega_P}
\newcommand{\sech}{\operatorname{sech}}

\title{V.F.S. v2.0 (Open-Gate)\\Lyapunov Core Proof Formulae}
\author{Core extraction from \texttt{lyapunov.html}}
\date{}

\begin{document}
\maketitle

\begin{abstract}
This document collects the core formulae of the Lyapunov proof for V.F.S. v2.0 (Open-Gate). It is not a separate proof text and contains no graphs. It records the base dynamics, normalized variables, Lyapunov functional, active-domain walls, radial corridor, bounded Open-Gate source, memory bounds, non-Zeno estimate, and asymptotic cleansing statements.
\end{abstract}

\tableofcontents
\newpage

% ============================================================
\section{Notation and Core Variables}

\begin{longtable}{>{\raggedright\arraybackslash}p{0.22\linewidth} >{\raggedright\arraybackslash}p{0.68\linewidth}}
\toprule
Symbol & Meaning \\
\midrule
$V$ & Voluntas: will coordinate. \\
$F$ & Factum: action coordinate. \\
$\sigma$ & Singular resistance. \\
$\lambda$ & Sophia: wisdom coordinate. \\
$\Om$ & Brim / dynamic capacity of Pleroma. \\
$u$ & Synergic intensity from $V$ and $F$. \\
$\Lambda_c$ & Transformation threshold. \\
$\mu$ & Saturation margin relative to the death-boundary. \\
$q$ & Normalized Sophia, $q=\lambda/\Om$. \\
$r$ & Normalized intensity, $r=u/\Om$. \\
$h_1,h_2$ & Normalized will and action. \\
$h_3$ & Normalized threshold variable. \\
\bottomrule
\end{longtable}

Core threshold and intensity:
\begin{equation}
  u=\sqrt{VF+\varepsilon},\qquad
  \Lambda_c=\frac{\gamma}{\delta},\qquad
  \varepsilon\ge 0.
\end{equation}

Transformation/collapse reading:
\begin{equation}
  u>\Lambda_c\quad\text{transformation region},
  \qquad
  u<\Lambda_c\quad\text{collapse region}.
\end{equation}

% ============================================================
\section{Base V.F.S. Dynamics}

Resistance equation:
\begin{equation}
  \dot\sigma=(\gamma-\delta u)\tanh(\kappa\sigma)
  =-\delta(u-\Lambda_c)\tanh(\kappa\sigma).
\end{equation}

Symmetric will--action dynamics:
\begin{equation}
  \dot V=\mu(aF+c\lambda)-\rho V,
  \qquad
  \dot F=\mu(aV+c\lambda)-\rho F,
  \qquad
  \mu=1-\frac{u}{\Om}.
\end{equation}

Capacity / Brim equation:
\begin{equation}
  \dot\Om=\alpha\lambda.
\end{equation}

Closed Microcosm limit:
\begin{equation}
  \dot\lambda=-\dot\sigma=
  \delta(u-\Lambda_c)\tanh(\kappa\sigma),
  \qquad
  \dot\sigma+
  \dot\lambda=0.
\end{equation}

Open-Gate Sophia dynamics:
\begin{equation}
  \dot\lambda=-\dot\sigma+
  \Igate(t)
  =\delta(u-\Lambda_c)\tanh(\kappa\sigma)+\Igate(t).
\end{equation}

Open balance:
\begin{equation}
  \dot\sigma+
  \dot\lambda=\Igate(t)\ge0,
  \qquad
  \sigma(t)+\lambda(t)=\sigma_0+\lambda_0+\Grec(t).
\end{equation}

Received grace integral:
\begin{equation}
  \Grec(t)=\int_0^t \Igate(s)\,ds,
  \qquad
  \dot{\Grec}=\Igate.
\end{equation}

% ============================================================
\section{Open-Gate Source}

Positive part of Sophia:
\begin{equation}
  \lambda_+=\frac{1}{m}\ln(1+e^{m\lambda}),
  \qquad m>0.
\end{equation}

Effective gate:
\begin{equation}
  g_{\rm eff}(t)=e^{-\tau t}+\left(1-e^{-\tau t}\right)\frac{H_K}{H_s+H_K}.
\end{equation}

Live Open-Gate source:
\begin{equation}
  \Igate(t)=\zeta_0\,g_{\rm eff}(t)e^{-\phi\sigma(t)}
  \frac{1}{1+\chi\lambda_+(t)}.
\end{equation}

Boundedness of the gate:
\begin{equation}
  H_K\ge0,\quad H_s>0
  \quad\Longrightarrow\quad
  0\le \frac{H_K}{H_s+H_K}<1,
  \qquad
  0\le g_{\rm eff}(t)\le1.
\end{equation}

Boundedness of live inflow:
\begin{equation}
  0\le \Igate(t)\le \zeta_0.
\end{equation}

% ============================================================
\section{Resistance Potential}

Resistance potential:
\begin{equation}
  \Psi_\sigma(\sigma)=\int_0^\sigma\tanh(\kappa s)\,ds
  =\frac{1}{\kappa}\ln\cosh(\kappa\sigma)\ge0.
\end{equation}

Derivative of the resistance potential:
\begin{equation}
  \dot\Psi_\sigma
  =\Psi_\sigma'(\sigma)\dot\sigma
  =\tanh(\kappa\sigma)(\gamma-\delta u)\tanh(\kappa\sigma)
  =(\gamma-\delta u)\tanh^2(\kappa\sigma).
\end{equation}

Active transformation margin:
\begin{equation}
  u\ge\Lambda_c+\eta,
  \qquad \eta>0.
\end{equation}

Dissipation in the active margin:
\begin{equation}
  \dot\Psi_\sigma\le -\delta\eta\tanh^2(\kappa\sigma).
\end{equation}

On $0\le\sigma\le C_\sigma$, there exists $c_\sigma>0$ such that:
\begin{equation}
  \dot\Psi_\sigma\le -c_\sigma\Psi_\sigma.
\end{equation}

Hence:
\begin{equation}
  \Psi_\sigma(t)\le \Psi_\sigma(0)e^{-c_\sigma t}.
\end{equation}

Near $\sigma=0$:
\begin{equation}
  \Psi_\sigma(\sigma)\sim \frac{\kappa}{2}\sigma^2,
  \qquad
  \sigma(t)=O\bigl(e^{-c_\sigma t/2}\bigr).
\end{equation}

% ============================================================
\section{Absolute Barrier Functional}

The absolute barrier is a motivational, non-primary barrier guarding absolute boundary faces:
\begin{equation}
\begin{aligned}
  \mathcal L_{abs}={}&
  A_\sigma\Psi_\sigma+E_VV^2+E_FF^2+E_\lambda\lambda^2
  +E_\Delta\frac{(V-F)^2}{2} \\
  &-\eta_V\ln V-\eta_F\ln F
  -\eta_\Omega\ln(\Om-\Lambda_c)
  -\eta_m\ln(\Om-u).
\end{aligned}
\end{equation}

Guarded absolute faces:
\begin{equation}
  V=0,
  \qquad
  F=0,
  \qquad
  \Om=\Lambda_c,
  \qquad
  u=\Om.
\end{equation}

% ============================================================
\section{Normalized Variables}

Normalized coordinates:
\begin{equation}
  h_1=\frac{V}{\Om},
  \qquad
  h_2=\frac{F}{\Om},
  \qquad
  q=\frac{\lambda}{\Om},
  \qquad
  r=\frac{u}{\Om},
  \qquad
  \mu=1-r,
  \qquad
  h_3=1-\frac{\Lambda_c}{\Om}.
\end{equation}

Product relation:
\begin{equation}
  r^2=h_1h_2+e(h_3),
  \qquad
  e(h_3)=\varepsilon\frac{(1-h_3)^2}{\Lambda_c^2}.
\end{equation}

Closed normalized dynamics:
\begin{equation}
  \dot h_1=\mu(ah_2+cq)-\rho h_1-\alpha qh_1,
\end{equation}
\begin{equation}
  \dot h_2=\mu(ah_1+cq)-\rho h_2-\alpha qh_2,
\end{equation}
\begin{equation}
  \dot q=\delta(r+h_3-1)\tanh(\kappa\sigma)-\alpha q^2,
\end{equation}
\begin{equation}
  \dot h_3=\alpha q(1-h_3).
\end{equation}

Open-Gate normalized $q$-source:
\begin{equation}
  \Sq=\frac{\Igate}{\Om}
  =\frac{\zeta_0g_{\rm eff}e^{-\phi\sigma}}
  {(1+\chi\lambda_+)\Om}.
\end{equation}

Open-Gate normalized $q$-equation:
\begin{equation}
  \dot q=\delta(r+h_3-1)\tanh(\kappa\sigma)-\alpha q^2+\Sq.
\end{equation}

Bound on $S_q$ inside the active domain:
\begin{equation}
  0\le \Sq\le \frac{\zeta_0}{\Om}
  \le \frac{\zeta_0(1-h_{3,*})}{\Lambda_c}
  =:S_q^{\max}.
\end{equation}

% ============================================================
\section{Product-Balance Mechanism}

Balance defect:
\begin{equation}
  D_\Delta=\frac12(h_1-h_2)^2.
\end{equation}

Let:
\begin{equation}
  d=h_1-h_2.
\end{equation}

Difference equation:
\begin{equation}
  \dot d=-(a\mu+\rho+\alpha q)d.
\end{equation}

Dissipation of balance defect:
\begin{equation}
  \dot D_\Delta=-2(a\mu+\rho+\alpha q)D_\Delta.
\end{equation}

If $\mu\ge\mu_*>0$ and $a>0$:
\begin{equation}
  \dot D_\Delta\le -2a\mu_*D_\Delta.
\end{equation}

Exponential will--action alignment:
\begin{equation}
  D_\Delta(t)\le D_\Delta(0)e^{-2a\mu_*t},
  \qquad
  |h_1(t)-h_2(t)|\le |h_1(0)-h_2(0)|e^{-a\mu_*t}.
\end{equation}

Product floor from radial and threshold walls:
\begin{equation}
  \Om\ge \Omega_{\min}:=\frac{\Lambda_c}{1-h_{3,*}}.
\end{equation}

If $r\ge r_*$ and $h_3\ge h_{3,*}$:
\begin{equation}
  h_1h_2=r^2-\frac{\varepsilon}{\Om^2}
  \ge p_*:=r_*^2-\frac{\varepsilon}{\Omega_{\min}^2}.
\end{equation}

Let:
\begin{equation}
  d_*:=\sqrt{2D^*}.
\end{equation}

Product-balance lower bound:
\begin{equation}
  h_1,h_2\ge m_{pb}:=\frac{\sqrt{d_*^2+4p_*}-d_*}{2}>0.
\end{equation}

% ============================================================
\section{Lyapunov Functional}

Product-balance Lyapunov core:
\begin{equation}
  \mathcal L_{pb}=A_\sigma\Psi_\sigma(\sigma)+A_\Delta D_\Delta+A_q q^2.
\end{equation}

Radial logarithmic barrier:
\begin{equation}
  B(r)=-\ln(1-r),
  \qquad
  0<r<1,
  \qquad
  r=\frac{u}{\Om}.
\end{equation}

Non-normalized radial barrier:
\begin{equation}
  B(u,\Om)=-\ln\left(1-\frac{u}{\Om}\right)
  =\ln\left(\frac{\Om}{\Om-u}\right).
\end{equation}

Primary V.F.S. Lyapunov functional:
\begin{equation}
  \boxed{
  \LV=\mathcal L_{pb}+A_rB(r)
  =A_\sigma\Psi_\sigma(\sigma)+A_\Delta D_\Delta+A_q q^2+A_rB(r),
  \qquad
  A_\sigma,A_\Delta,A_q,A_r>0.}
\end{equation}

Expanded form:
\begin{equation}
  \LV=A_\sigma\Psi_\sigma(\sigma)+A_\Delta D_\Delta+A_q q^2
  -A_r\ln(1-r).
\end{equation}

Non-normalized expanded form:
\begin{equation}
  \LV=A_\sigma\Psi_\sigma(\sigma)+A_\Delta D_\Delta+A_q q^2
  -A_r\ln\left(1-\frac{u}{\Om}\right).
\end{equation}

Radial boundary protection:
\begin{equation}
  B(r)\to+\infty\quad\text{as}\quad r\to1^-.
\end{equation}

% ============================================================
\section{Active Domain}

Active domain:
\begin{equation}
\DA=\left\{
\begin{array}{l}
0\le\sigma\le C_\sigma,\quad q\in[0,q^*],\quad h_3\in[h_{3,*},1),\\
r\in[r_*,r^*],\quad \mu=1-r\ge\mu_*>0,\\
D_\Delta\le D^*,\quad h_1>0,\quad h_2>0
\end{array}
\right\},
\qquad
0<r_*<r^*<1.
\end{equation}

Equivalent death-boundary avoidance inside $\DA$:
\begin{equation}
  0<r<1
  \qquad\Longleftrightarrow\qquad
  0<u<\Om.
\end{equation}

Closed upper $q$-wall:
\begin{equation}
  \alpha(q^*)^2\ge\delta r^*.
\end{equation}

Open-Gate upper $q$-wall:
\begin{equation}
  \boxed{\alpha(q^*)^2\ge \delta r^*+S_q^{\max}.}
\end{equation}

Lower $q$-wall in Open-Gate:
\begin{equation}
  q=0:
  \qquad
  \dot q=\delta(r+h_3-1)\tanh(\kappa\sigma)+S_q\ge0
\end{equation}
provided the active margin makes $r+h_3-1\ge0$.

% ============================================================
\section{Radial Corridor}

Radial identity from $r^2=h_1h_2+e(h_3)$:
\begin{equation}
  2r\dot r=
  \mu\bigl[a(h_1^2+h_2^2)+cq(h_1+h_2)\bigr]
  -2\rho(r^2-e(h_3))-2\alpha q r^2
  =:\Rcal_{\rm rad}.
\end{equation}

Maximum $e$ on the active corridor:
\begin{equation}
  e_*^{\max}=\varepsilon\frac{(1-h_{3,*})^2}{\Lambda_c^2}.
\end{equation}

Non-degeneracy:
\begin{equation}
  r_*^2>e_*^{\max}.
\end{equation}

Active margin:
\begin{equation}
  r_*+h_{3,*}-1\ge \eta\frac{1-h_{3,*}}{\Lambda_c}.
\end{equation}

Lower radial wall:
\begin{equation}
  a(1-r_*)\bigl(r_*^2-e_*^{\max}\bigr)
  \ge (\rho+\alpha q^*)r_*^2.
\end{equation}

This wall is sufficient, not tight. The exact lower-wall requirement is
\begin{equation}
  (1-r_*)a(r_*^2-e)
  \ge
  \rho(r_*^2-e)+\alpha q r_*^2,
\end{equation}
and the stated form replaces \(\rho(r_*^2-e)\) by the larger \(\rho r_*^2\).

Upper radial wall:
\begin{equation}
  (1-r^*)\bigl[a(d_*^2+2r^{*2})+cq^*\sqrt{d_*^2+4r^{*2}}\bigr]
  \le 2\rho\bigl(r^{*2}-e_*^{\max}\bigr).
\end{equation}

Direct upper radial inwardness condition:
\begin{equation}
  \mu\left[a(h_1^2+h_2^2)+cq(h_1+h_2)\right]
  \le 2\rho(r^{*2}-e)+2\alpha q r^{*2}
  \qquad\text{on } r=r^*.
\end{equation}

Open-Gate radial convention:
\begin{equation}
  q_{\rm new}^*:=\sqrt{\frac{\delta r^*+S_q^{\max}}{\alpha}}.
\end{equation}

Conservative lower radial wall with $q_{\rm new}^*$:
\begin{equation}
  a(1-r_*)\bigl(r_*^2-e_*^{\max}\bigr)
  \ge (\rho+\alpha q_{\rm new}^*)r_*^2.
\end{equation}

Epektasis-dominance derivative at the upper radial wall:
\begin{equation}
  \frac{\partial\Rcal}{\partial q}\bigg|_{r^*}
  =\mu c(h_1+h_2)-2\alpha r^{*2}.
\end{equation}

Worst-case bound:
\begin{equation}
  \frac{\partial\Rcal}{\partial q}\bigg|_{r^*}
  \le (1-r^*)c\sqrt{d_*^2+4r^{*2}}-2\alpha r^{*2}.
\end{equation}

Epektasis-dominance sufficient condition:
\begin{equation}
  \boxed{
  2\alpha r^{*2}\ge (1-r^*)c\sqrt{d_*^2+4r^{*2}}
  \quad\Longrightarrow\quad
  \frac{\partial\Rcal}{\partial q}\le0.}
\end{equation}

% ============================================================
\section{Lyapunov Dissipative Estimate}

Derivative decomposition:
\begin{equation}
\begin{aligned}
  \dot\LV={}&
  A_\sigma\tanh(\kappa\sigma)\dot\sigma
  +A_\Delta\dot D_\Delta
  +2A_qq\dot q
  +A_r\dot B(r).
\end{aligned}
\end{equation}

Open-Gate contribution to the $q^2$ block:
\begin{equation}
  \frac{d}{dt}(A_q q^2)
  =2A_q q\left[\delta(r+h_3-1)\tanh(\kappa\sigma)-\alpha q^2\right]
  +2A_q qS_q.
\end{equation}

Bound on Open-Gate contribution:
\begin{equation}
  0\le 2A_q qS_q\le 2A_q q^*S_q^{\max}.
\end{equation}

Dissipative Lyapunov estimate:
\begin{equation}
  \boxed{
  \dot\LV\le C-C_1\LV,
  \qquad
  C_1>0,
  \qquad
  C<\infty.}
\end{equation}

Open-Gate constant split:
\begin{equation}
  C=C_{\rm closed}+\Delta C_{\rm gate},
  \qquad
  0\le\Delta C_{\rm gate}\le 2A_q q^*S_q^{\max}.
\end{equation}

Key structural point:
Here $C_1$ is unchanged by the Open-Gate source, while $C$ receives only a bounded gate contribution.

% ============================================================
\section{Main Theorem Formulae}

Positive invariance:
\begin{equation}
  x(0)\in\DA
  \quad\Longrightarrow\quad
  x(t)\in\DA\quad\forall t\ge0.
\end{equation}

Reset admissibility holds by the Reset Lemma, with
\begin{equation}
  L_R=\sup_{\DA}\LV.
\end{equation}

Bound from the Lyapunov inequality:
\begin{equation}
  \LV(t)\le L_*:=\max\left\{L_R,\frac{C}{C_1}\right\}.
\end{equation}

Radial barrier implication:
\begin{equation}
  A_rB(r)\le L_*.
\end{equation}

Explicit radial bound:
\begin{equation}
  -\ln(1-r)\le\frac{L_*}{A_r}
  \quad\Longrightarrow\quad
  r(t)\le1-e^{-L_*/A_r}<1.
\end{equation}

Death-boundary avoidance:
\begin{equation}
  r(t)<1
  \quad\Longrightarrow\quad
  u(t)<\Om(t).
\end{equation}

Product-balance positivity:
\begin{equation}
  h_1,h_2\ge m_{pb}>0.
\end{equation}

No artificial $q$-floor:
\begin{equation}
  q\ge0\quad\text{suffices; no assumption }q\ge q_0>0\text{ is required.}
\end{equation}

% ============================================================
\section{Reset Admissibility (Resurrectio Jump)}

The resets $\Rcal$ are \emph{exogenous}: since $\DA$ is positively invariant, the
continuous flow never reaches the boundary of the active domain by itself. A death-like
strike is therefore modeled as occurring at an interior point
\begin{equation}
  x^-=(\sigma^-,\lambda^-,V,F,\Om^-)\in\DA
\end{equation}
at an unpredictable time.

\textbf{Reset map (linear metabolism).} With $0\le q_R<1$ and $\kappa_R>0$,
\begin{equation}
  \Rcal:\quad
  \sigma^+=q_R\sigma^-,\qquad
  \lambda^+=\lambda^-+(1-q_R)\sigma^-,\qquad
  (V,F)\ \text{continuous},\qquad
  \Om^+=\Om^-+\kappa_R\Grec^-.
\end{equation}
Here
\begin{equation}
  \Grec^-:=\Grec(t_{\rm death}^-),
  \qquad
  u^+=\sqrt{VF+\varepsilon}=u^-.
\end{equation}

\textbf{Quantities.}
\begin{equation}
  A_R:=\Grec^-,\qquad
  R_c:=C_\sigma-\sigma_0-\lambda_0,
\end{equation}
\begin{equation}
  R_{\min}^{\rm ceil}:=
  \frac{(1-q_R)\sigma^-}{q^*\kappa_R},
  \qquad
  R_{\max}:=
  \frac{\Om^-(r^--r_*)}{\kappa_R r_*}.
\end{equation}
By the Open-Gate balance
\begin{equation}
  \sigma^-+\lambda^-=
  \sigma_0+\lambda_0+\Grec^-,
\end{equation}
one has
\begin{equation}
  \Grec^-\ge R_c
  \quad\Longleftrightarrow\quad
  \sigma^-+\lambda^-\ge C_\sigma
  \quad\Longleftrightarrow\quad
  \lambda^-\ge C_\sigma-\sigma^-.
\end{equation}

\textbf{Lemma (reset admissibility).}
Let $\kappa_R>0$ and $x^-\in\DA$. If
\begin{equation}
  \max\{R_c,R_{\min}^{\rm ceil}\}\le\Grec^-\le R_{\max},
\end{equation}
then $x^+:=\Rcal(x^-)\in\DA$ and
\begin{equation}
  \LV(x^+)\le L_R:=\sup_{\DA}\LV<\infty.
\end{equation}

\textbf{Proof.} Coordinate by coordinate,
\begin{equation}
  \sigma^+=q_R\sigma^-\in[0,C_\sigma).
\end{equation}
Moreover,
\begin{equation}
  h_3^+=1-\frac{\Lambda_c}{\Om^+}
  \ge h_3^-\ge h_{3,*},
  \qquad
  h_3^+<1,
\end{equation}
since $\Om^+\ge\Om^-$. Also
\begin{equation}
  h_1^+=h_1^-\frac{\Om^-}{\Om^+},
  \qquad
  h_2^+=h_2^-\frac{\Om^-}{\Om^+},
\end{equation}
so
\begin{equation}
  D_\Delta^+=\frac12\left(d^-\frac{\Om^-}{\Om^+}\right)^2
  \le D_\Delta^-\le D^*,
  \qquad
  h_1^+,h_2^+>0.
\end{equation}
The radial coordinate satisfies
\begin{equation}
  r^+=\frac{r^-\Om^-}{\Om^-+\kappa_R\Grec^-}
  \le r^-\le r^*.
\end{equation}
The lower radial wall is preserved exactly when
\begin{equation}
  r^+\ge r_*
  \quad\Longleftrightarrow\quad
  \kappa_R\Grec^-\le\frac{\Om^-(r^--r_*)}{r_*}
  \quad\Longleftrightarrow\quad
  \Grec^-\le R_{\max}.
\end{equation}
Finally,
\begin{equation}
  q^+=
  \frac{\lambda^-+(1-q_R)\sigma^-}{\Om^-+\kappa_R\Grec^-}
  \ge0.
\end{equation}
Since $x^-\in\DA$ gives $\lambda^-\le q^*\Om^-$,
\begin{equation}
\begin{aligned}
  q^*(\Om^-+\kappa_R\Grec^-)
  -\bigl(\lambda^-+(1-q_R)\sigma^-\bigr)
  &\ge
  q^*\kappa_R\Grec^--(1-q_R)\sigma^-\\
  &\ge0
\end{aligned}
\end{equation}
whenever $\Grec^-\ge R_{\min}^{\rm ceil}$. Hence $q^+\le q^*$.
Thus $x^+\in\DA$. The functional $\LV$ is bounded on $\DA$ because
\begin{equation}
  \Psi_\sigma\le\Psi_\sigma(C_\sigma),
  \qquad
  D_\Delta\le D^*,
  \qquad
  q^2\le(q^*)^2,
  \qquad
  B(r)\le-\ln(1-r^*)<\infty.
\end{equation}
Therefore $\LV(x^+)\le L_R$. \hfill$\square$

\textbf{Consequences.}
\begin{enumerate}[label=(\arabic*)]
\item The admissible window floor is
\begin{equation}
  \max\{R_c,R_{\min}^{\rm ceil}\}.
\end{equation}
It combines the theological gate and the geometric ceiling-floor: grace must expand the vessel sufficiently to absorb the converted resistance. For cleansed deaths, $\sigma^-\to0$, the ceiling-floor vanishes and $R_c$ alone remains binding.

\item The condition $\kappa_R>0$ suffices. A natural design choice is
\begin{equation}
  \kappa_R\ge\frac1{q^*},
\end{equation}
so that the asymptotic ceiling $q^+\to1/\kappa_R\le q^*$ and, in the intended window, $R_{\min}^{\rm ceil}\le R_c$.

\item The admissible grace window is
\begin{equation}
  \Grec^-\in
  \left[\max\{R_c,R_{\min}^{\rm ceil}\},R_{\max}\right].
\end{equation}
Below this window one obtains final death; above it one obtains over-dilution, $r<r_*$. The window is nonempty if and only if
\begin{equation}
  \Om^-(r^--r_*)
  \ge
  \kappa_R r_*\max\{R_c,R_{\min}^{\rm ceil}\}.
\end{equation}
A death pinned at $r^-=r_*$ is not resurrectible.

\item With the non-Zeno dwell time
\begin{equation}
  \Delta t_j\ge\frac{\bar h_*}{\bar M_*}>0,
\end{equation}
every life-arc satisfies
\begin{equation}
  \LV\le L_*=
  \max\left\{L_R,\frac{C}{C_1}\right\}=L_R,
\end{equation}
and resets do not accumulate. Hence the hybrid trajectory is forward complete across any infinite sequence of life-arcs.
\end{enumerate}

\textbf{Theological gloss.} This remark is not used in the proof. The gate
\begin{equation}
  \Grec^-\ge R_c=C_\sigma-\sigma_0-\lambda_0
\end{equation}
is the threshold past which Sophia can no longer be extinguished by resistance; in that regime $\lambda^->0$ is forced by the balance. Its complement, $\lambda^-\to0$, is the self-closing limit $\Igate\to0$: final impenitence. It resembles the oil of the wise virgins and the \emph{donum perseverantiae}. The upper boundary $R_{\max}$ represents the opposite failure: an over-diluting or ``cheap grace'' regime.

% ============================================================
\section{Non-Zeno Estimate}

Hybrid margin:
\begin{equation}
  \bar h_*=
  \min\left\{m_{pb},h_{3,*},r_*,1-r^*,q^*,C_\sigma,\sqrt{D^*}\right\}>0.
\end{equation}

Bounded vector field:
\begin{equation}
  \bar M_*<\infty.
\end{equation}

Dwell-time lower bound:
\begin{equation}
  \Delta t_j\ge \Delta t_*:=\frac{\bar h_*}{\bar M_*}>0.
\end{equation}

Non-Zeno conclusion:
\begin{equation}
  \sum_j \Delta t_j=+\infty.
\end{equation}

% ============================================================
\section{Memory and Gate Bounds}

Mensura kernel:
\begin{equation}
  R(z)=\frac{z}{z+K},
  \qquad
  z\ge0,
  \qquad
  0\le R(z)\le1.
\end{equation}

Saturated gate memory:
\begin{equation}
  \frac{H_K}{H_s+H_K}\in[0,1)
  \qquad
  \forall H_K\ge0.
\end{equation}

Gate saturation:
\begin{equation}
  g_{\rm eff}=e^{-\tau t}+(1-e^{-\tau t})\frac{H_K}{H_s+H_K}\in[0,1].
\end{equation}

Memory velocity bounds:
\begin{equation}
  |\dot H_K|\le \eta\left[R((\dot x_0)_+)+R((\dot x_0)_-)\right]\le 2\eta.
\end{equation}

Wound-memory velocity bound:
\begin{equation}
  |\dot W|\le2\eta.
\end{equation}

Bounded-gain loop:
\begin{equation}
  H_K\xrightarrow{\le1}g_{\rm eff}
  \xrightarrow{\le S_q^{\max}}S_q
  \to\lambda\to V,F\to u\to\dot x_0
  \xrightarrow{R\le1}\dot H_K.
\end{equation}

Extended field bound:
\begin{equation}
  \bar M_*<\infty
  \quad\text{including}\quad
  |\dot H_K|,|\dot W|\le2\eta.
\end{equation}

% ============================================================
\section{Asymptotic Cleansing and Alignment}

Will--action alignment:
\begin{equation}
  D_\Delta(t)\le D_\Delta(0)e^{-2a\mu_*t},
  \qquad
  |h_1(t)-h_2(t)|\le |h_1(0)-h_2(0)|e^{-a\mu_*t}\to0.
\end{equation}

Resistance cleansing under persistent active margin:
\begin{equation}
  u(t)\ge\Lambda_c+\eta\quad\forall t\ge0
  \quad\Longrightarrow\quad
  \Psi_\sigma(t)\le\Psi_\sigma(0)e^{-c_\sigma t}
  \quad\Longrightarrow\quad
  \sigma(t)\to0.
\end{equation}

\textbf{Corollary (unconditional cleansing inside the active domain).}
Since \(\DA\) is positively invariant, one has \(r(t)\ge r_*\) and
\(h_3(t)\ge h_{3,*}\) for all \(t\ge0\). Thus
\begin{equation}
  r(t)+h_3(t)-1
  \ge
  r_*+h_{3,*}-1
  \ge
  \frac{\eta(1-h_{3,*})}{\Lambda_c}>0.
\end{equation}
Equivalently,
\begin{equation}
  u(t)\ge\Lambda_c+\eta
  \qquad\forall t\ge0.
\end{equation}
Therefore the hypothesis of the persistent-margin cleansing statement holds automatically for every trajectory with \(x(0)\in\DA\). Hence
\begin{equation}
  \sigma(t)\to0,
  \qquad
  h_1(t)-h_2(t)\to0.
\end{equation}

Open-balance asymptotic Sophia:
\begin{equation}
  \lambda(t)=\sigma_0+\lambda_0+\Grec(t)-\sigma(t).
\end{equation}

If $\Grec(\infty)<\infty$ and $\sigma(t)\to0$:
\begin{equation}
  \lambda(t)\to \sigma_0+\lambda_0+\Grec(\infty).
\end{equation}

If $\liminf_{t\to\infty}\lambda(t)>0$:
\begin{equation}
  \Om(t)\to+\infty.
\end{equation}

Epektatic dilution of $S_q$:
\begin{equation}
  \int_0^\infty\lambda(t)\,dt=+\infty
  \quad\Longrightarrow\quad
  \Om(t)\to+\infty
  \quad\Longrightarrow\quad
  S_q(t)\to0.
\end{equation}

Proof of the last implication:
\begin{equation}
  0\le S_q(t)\le\frac{\zeta_0}{\Om(t)}\to0.
\end{equation}

% ============================================================

\section{Numerical Witness Formulae}

One admissible parameter point recorded in the Lyapunov file:
\begin{equation}
\begin{aligned}
&a=3.494,\quad c=0.429,\quad \alpha=0.529,\quad \delta=0.898,\quad \rho=0.545,\\
&\zeta_0=0.560,\quad h_{3,*}=0.622,\quad r_*=r_-=0.531,\quad r^*=r_+=0.892,\\
&\varepsilon=0.006,\quad d=0.193,\qquad \Lambda_c=1\quad\text{(normalization)}.
\end{aligned}
\end{equation}

Derived quantities:
\begin{equation}
\begin{aligned}
  e_{\max}&=\varepsilon(1-h_{3,*})^2=0.0009,\\
  \Sq^{\max}&=\zeta_0(1-h_{3,*})=0.212,\\
  q_{
m new}&=\sqrt{\frac{\delta r_++\Sq^{\max}}{\alpha}}=1.383,\\
  q^*&=1.05q_{
m new}=1.452.
\end{aligned}
\end{equation}

Active margin:
\begin{equation}
  r_-+h_{3,*}-1=0.153>0,
  \qquad
  \frac{u}{\Lambda_c}\approx1.40.
\end{equation}
Thus the witness lies in the Transformation region.

All five parametric walls hold strictly; the worst-case slack is
\begin{equation}
  0.084.
\end{equation}

% ============================================================
\section{Final Proof Chain}

Core proof chain:
\begin{equation}
\boxed{
\mathrm{Admissibility}
\Longrightarrow
\DA
\Longrightarrow
\dot\LV\le C-C_1\LV
\Longrightarrow
r(t)<1
\Longrightarrow
\mathrm{NonZeno}.}
\end{equation}

Reset admissibility holds by the Reset Lemma, with
\begin{equation}
  L_R=\sup_{\DA}\LV.
\end{equation}

Active transformation conclusion:
\begin{equation}
  u\ge\Lambda_c+\eta
  \quad\Longrightarrow\quad
  \sigma(t)\to0,
  \qquad
  h_1(t)-h_2(t)\to0.
\end{equation}

The theorem is an active-domain theorem, not global stability on all phase space:
The active-domain assumption is essential:
\begin{equation}
  x(0)\in\DA.
\end{equation}

Closed limit:
\begin{equation}
  \zeta_0=0
  \quad\Longrightarrow\quad
  S_q=0,
  \quad
  \Igate=0,
  \quad
  \dot\lambda=-\dot\sigma,
  \quad
  \sigma+
  \lambda=\sigma_0+
  \lambda_0.
\end{equation}

Open-Gate theorem summary:
Open-Gate adds only bounded $S_q$ in $\dot q$; it strengthens walls but preserves the dissipative core.

\end{document}
